\section{Merge: Pinned Guards and Open Trace Variables}\label{sec:merge}

\subsection{The program}
\begin{lstlisting}
method main(a: array<int, int>, m: int, b: array<int, int>, n: int) returns (c: array<int, int>)
  requires 0 <= m && 0 <= n
  requires sorted(a, 0, m) && sorted(b, 0, n)
  ensures sorted(c, 0, m + n)
  ensures forall x: int :: count(x, selects(c, butlast(counter(0, 1, natural(m + n)))))
    == count(x, selects(a, butlast(counter(0, 1, natural(m)))))
     + count(x, selects(b, butlast(counter(0, 1, natural(n)))))
{
  var i: int; var j: int; var k: int;
  i := 0; j := 0; k := 0;
  while k < m + n {
    if n <= j || (i < m && a[i] <= b[j]) { c := c[k := a[i]]; i := i + 1; }
    else { c := c[k := b[j]]; j := j + 1; }
    k := k + 1;
  }
}
\end{lstlisting}
The second postcondition states the permutation as a multiset: for each value, its number of occurrences
among the first $m + n$ elements of $c$ is the sum of its occurrences in $a$ and $b$, in the vocabulary of
\texttt{teq.dfy}.

\subsection{The trace equation system}
\begin{align}
  K &= \cnt(0, 1, n') \qquad c = \cur(C) \qquad k = \intf(n') \qquad m + n \le k \label{eq:mg0}\\
  C &= \tr{c_0}{\bl{K}, \bl{I}, \bl{J}}{c, k, i, j \mapsto \upd{G}{c[k := a[i]]}\ \ \upd{\neg G}{c[k := b[j]]}} \label{eq:mg1}\\
  I &= \tr{0}{\bl{K}, \bl{J}}{i, k, j \mapsto \upd{G}{i + 1}\ \ \upd{\neg G}{i}} \label{eq:mg2}\\
  J &= \tr{0}{\bl{K}, \bl{I}}{j, k, i \mapsto \upd{G}{j}\ \ \upd{\neg G}{j + 1}} \label{eq:mg3}
\end{align}
where $G$ abbreviates the test $n \le j \vee (i < m \wedge a[i] \le b[j])$ of the program. The pointers
are traces over each other, \eqref{eq:mg2} and \eqref{eq:mg3}: $I$ has $\bl{J}$ among its dependencies and $J$ has $\bl{I}$. A single loop,
no fragments; the whole difficulty is in matching this shape.

\subsection{The lemmas}\label{sec:merge:lemmas}
The merge is named by an abstract list $\mathsf{merged}(a, m, b, n)$, the values of the result on
$0, \ldots, m + n - 1$, with two facts: it is sorted, and its counts are those of both ranges. The
lemma on the loop rewrites the values of the result:
\begin{align*}
  &\lemma{selects\_merge\_loop}:\\
  &\quad s = m + n\ \wedge\ m + n \le \intf(nn)\ \wedge\ 0 \le m \wedge 0 \le n\\
  &\quad\wedge\ (\forall i, j.\ G(i, j) \Leftrightarrow (n \le j \vee (i < m \wedge a[i] \le b[j])))\\
  &\quad\wedge\ (\forall i, j.\ U(i, j) = a[i])\ \wedge\ (\forall i, j.\ V(i, j) = b[j])\\
  &\quad\wedge\ Iv = \tr{0}{\bl{\cnt(0, 1, nn)}, js}{i, k, j \mapsto \upd{G(i, j)}{i + 1}\ \ \upd{\neg G(i, j)}{i}}\\
  &\quad\wedge\ is = \bl{Iv}\\
  &\quad\wedge\ Jv = \tr{0}{\bl{\cnt(0, 1, nn)}, is}{j, k, i \mapsto \upd{G(i, j)}{j}\ \ \upd{\neg G(i, j)}{j + 1}}\\
  &\quad\wedge\ js = \bl{Jv}\\
  &\quad\Longrightarrow\ \mathsf{selects}(\cur(\tr{c_0}{\bl{\cnt(0, 1, nn)}, is, js}{c, k, i, j \mapsto {}\\
  &\qquad\qquad\qquad \upd{G(i, j)}{c[k := U(i, j)]}\ \ \upd{\neg G(i, j)}{c[k := V(i, j)]}}), \bl{\cnt(0, 1, \natf(s))})\\
  &\qquad\qquad = \mathsf{merged}(a, m, b, n)
\end{align*}
and \lemma{sorted\_merge\_loop} is the same with the conclusion $\sorted(\cur(\ldots), 0, s) \Leftrightarrow \mathsf{true}$
and the two $\sorted$ premises; \lemma{count\_merged} gives
\begin{align*}
  \mathsf{count}(x, \mathsf{merged}(a, m, b, n)) = {}&\mathsf{count}(x, \mathsf{selects}(a, \bl{\cnt(0, 1, \natf(m))}))\\
  &+ \mathsf{count}(x, \mathsf{selects}(b, \bl{\cnt(0, 1, \natf(n))})).
\end{align*}

Three workarounds of the matcher are built into this lemma, each isolated by experiments
(\texttt{NOTES.md}, findings 14, 16 and the section on the pin):
\begin{enumerate}[nosep]
\item \textbf{The pin.} A literal disjunctive guard $n \le j \vee (\ldots)$ in a pattern is never
  matched. The guard is the higher-order variable $G(i, j)$, which matches any guard over the bound
  variables $i, j$, and the quantified side condition
  $\forall i, j.\ G(i, j) \Leftrightarrow (n \le j \vee \ldots)$ pins it to the intended one; the simplifier
  decides the pin by beta reduction. The stored values $U(i, j) = a[i]$ and $V(i, j) = b[j]$ are pinned
  likewise. (A higher-order variable for an array read $a[i]$ alone would be instantiated by the array
  itself, finding 10, hence the pin instead of a literal $a[i]$.)
\item \textbf{The open trace variable.} The lemma must know that $is$ and $js$ are the pointer traces,
  which the pattern cannot say, since each is defined over the other. A closed side condition
  $is = \bl{\tr{0}{\ldots}{\ldots}}$ is decided by the simplifier and fails (it does not expand $\bl{I}$).
  Instead $Iv = \tr{0}{\ldots, js}{\ldots}$ with an \emph{open} trace variable $Iv$ is matched against
  the hypothesis $I = \tr{0}{\bl{K}, \bl{J}}{\ldots}$ \eqref{eq:mg2}, which binds $Iv := I$ (and the
  counter $\bl{\cnt(0, 1, nn)}$ matches $\bl{K}$ through its definition); then $is = \bl{Iv}$ is closed
  and trivially true.
\item \textbf{The whole atom as the left-hand side.} A variable argument of an uninterpreted function
  is not expanded by its definition when the simplifier rewrites below it: $\sorted(c, 0, m + n)$ with
  $c = \cur(C)$ is not matched by a left-hand side $\cur(\tr{\ldots}{\ldots})$, but
  $\sorted(\cur(\tr{\ldots}{\ldots}), 0, s)$ as the whole atom is, and $\mathsf{selects}(\cur(\ldots), \ldots)$
  with $\mathsf{selects}$ as the head, inside $\mathsf{count}$ and the quantifier. The length is a
  variable $s$ with $s = m + n$ as a side condition, since a pattern $m + n$ of two variables is not matched.
\end{enumerate}
The remaining condition $m + n \le \intf(nn)$ is the exit $m + n \le k$ with $k = \intf(n')$ of
\eqref{eq:mg0}: after at least $m + n$ iterations the positions $0, \ldots, m + n - 1$ are those of the
merge, whatever further iterations store beyond them, which makes the lemma sound without fixing the
number of iterations exactly (the goal knows it only as a disjunction).

\subsection{The idea}
The loop computes a function of its inputs that has a name in the specification's vocabulary,
$\mathsf{merged}$, and the lemma database relates the loop to that name and the name to the properties.
The matching problems of mutually dependent traces are solved by two general devices, the pin and the
open trace variable, which apply to every literal pattern that the matcher does not handle: in
\texttt{insertion\_sort.dfy} with one pointer, the positions $j$ and $j - 1$ of the swaps are pinned in
the same way, which closed a file that had been open since the start.
